\section{Discussion}

Re-prompt Tax changes the unit of analysis from a single answer to a repair
trajectory. This matters because a first answer can be semantically plausible
but interactionally wrong. For global users, the cost of that error includes
not only extra tokens, but also the cognitive work of diagnosing the failure,
writing a corrective prompt, and verifying the repair. These costs are easy to
miss when benchmarks ask only whether a final response is acceptable.

The strongest empirical finding is that implicit content-language preservation
is a measurable bottleneck. Models often follow the language of the instruction
rather than the language of the quoted content to be edited. A simple contract
prompt reduces this failure, but does not eliminate exact-preservation errors.
This pattern suggests that multilingual system prompts should separate
\emph{instruction language}, \emph{content language}, and \emph{response
language} as distinct variables rather than treating the prompt as a monolithic
language signal.

The results also show why mitigation claims should be bounded. The full Global
Interaction Contract improves all three models, but a narrower content-
preservation prompt outperforms it on the nano diagnostic. This suggests that
more specific scaffolds may be better than broad multilingual policy text for
some workflows. Future systems could infer a lightweight interaction contract
from the user's prompt and attach only the relevant constraints, reducing both
absolute prompt length and over-instruction risk.
